% =====================================================================
% Clase -- Física 11° -- Trimestre I -- Semana 03
% Sesiones 004 (lun 14 sep., 11:00, 100 min) y 005 (mié 16 sep., 10:10, 100 min)
% Tema 2 de la guía: La carga eléctrica · hilo «De Tales al pararrayos»
% Material del docente (no se entrega a los estudiantes).
% =====================================================================
\documentclass[11pt]{article}
\makeatletter\def\input@path{{../../../../../plantillas/estilo/}}\makeatother
\usepackage{mmcantillo}
\guiadelestudiante{../../guia-didactica/guia-periodo-I-electrostatica}

\tipodocumento{Clase}
\asignatura{Física}
\titulo{Semana 3: la carga eléctrica y cómo se electriza un cuerpo}
\grado{11°}
\periodo{I}
% DBA: naturales grado 11 · DBA 2 — Comprende que la interacción de las cargas en reposo
%      genera fuerzas eléctricas y que cuando las cargas están en movimiento genera fuerzas
%      magnéticas.
% Evidencia trabajada: (1) identifica el tipo de carga eléctrica (positiva o negativa) que
%      adquiere un material cuando se somete a procedimientos de fricción o contacto.

\begin{document}

\encabezado*

\begin{center}
{\renewcommand{\arraystretch}{1.3}
\begin{tabularx}{\linewidth}{@{}l l l X l@{}}
  \toprule
  \textbf{Sesión} & \textbf{Fecha} & \textbf{Min} & \textbf{Subtema} & \textbf{Entregables}\\
  \midrule
  004 & lun 14 sep., 11:00 & 100 & Naturaleza de la carga, tipos y conservación; el ámbar de Tales & Taller 3\\
  005 & mié 16 sep., 10:10 & 100 & Electrización por frotamiento y por contacto; conductores y aislantes & Tarea 3 (para el lun 21 sep.)\\
  \bottomrule
\end{tabularx}}
\end{center}

Referencia: \refguia{tema:carga}; el repaso de notación científica usa la
\refguia{tema:notacion}.

% ---------------------------------------------------------------------
\seccion{Sesión 004 — lunes 14 de septiembre · 11:00 · 100 min}

\textbf{Objetivo.} Reconocer que hay dos clases de carga, que la carga se conserva y que está
cuantizada ($q = n\,e$), y calcular la carga de un cuerpo a partir de los electrones que gana
o pierde.

\begin{center}
{\renewcommand{\arraystretch}{1.3}
\begin{tabularx}{\linewidth}{@{}l l X@{}}
  \toprule
  \textbf{Momento} & \textbf{Min} & \textbf{Actividad}\\
  \midrule
  Apertura & 0–10 & Hilo: un tubo de PVC (o un globo) frotado con lana atrae papelitos, como
    el ámbar de Tales. En griego el ámbar se llamaba \emph{élektron}: de ahí «electricidad» y
    «electrón». Aclara que lo de Tales es una tradición (nos llega por Diógenes Laercio), no un
    dato seguro. Pregunta al grupo: ¿qué le pasó al tubo?\\
  Explicación & 10–35 & Dos clases de carga y ley de signos (demostración: dos tiras de cinta
    adhesiva despegadas de la mesa se repelen; una despegada de encima de la otra, se atraen).
    El átomo: protones en el núcleo, electrones alrededor; un cuerpo se carga ganando o
    perdiendo \emph{electrones}, nunca protones. Unidad: coulomb. Conservación (lo que gana uno
    lo pierde otro) y cuantización, $q = n\,e$ con $e = \qty{1,602e-19}{C}$ (en los
    ejercicios se redondea a $\qty{1,6e-19}{C}$).\\
  Ejemplo & 35–45 & Resuelve con el grupo el ejemplo «El globo y la lana» de la guía: el
    globo gana $\num{2e6}$ electrones y queda con $\qty{-3,2e-13}{C}$; la lana, con
    $\qty{+3,2e-13}{C}$. Insiste en el signo y en que las dos cargas suman cero.\\
  Repaso & 45–55 & Notación científica con cargas: ¿cuántos coulombs es un
    nanocoulomb ($\qty{1e-9}{C}$) y un microcoulomb ($\qty{1e-6}{C}$)? Recuerda la
    multiplicación de potencias de 10 de la semana 2.\\
  Taller 3 & 55–90 & En parejas.\\
  Cierre & 90–100 & Punto 4 del taller en el tablero: ¿por qué no puede existir una carga de
    $\qty{2,4e-19}{C}$? Anuncia que el miércoles se frotan materiales para ver qué signo toma
    cada uno.\\
  \bottomrule
\end{tabularx}}
\end{center}

\begin{nota}
Lleva un tubo de PVC (o un globo), un trozo de lana, papelitos y cinta adhesiva. La
demostración de las cintas: pega dos tiras en la mesa, despégalas rápido y acércalas (se
repelen); luego pega una sobre otra, despégalas entre sí y acércalas (se atraen).
\end{nota}

% ---------------------------------------------------------------------
\seccion{Sesión 005 — miércoles 16 de septiembre · 10:10 · 100 min}

\textbf{Objetivo.} Predecir el signo que adquiere un material al frotarlo con otro (serie
triboeléctrica) y al tocar un cuerpo cargado, y distinguir conductores de aislantes.

\begin{center}
{\renewcommand{\arraystretch}{1.3}
\begin{tabularx}{\linewidth}{@{}l l X@{}}
  \toprule
  \textbf{Momento} & \textbf{Min} & \textbf{Actividad}\\
  \midrule
  Repaso & 0–10 & Dos preguntas rápidas: ¿qué gana o pierde un cuerpo que queda positivo?
    ¿Cuántos electrones hay en $\qty{-3,2e-19}{C}$? (dos, en exceso).\\
  Frotamiento & 10–35 & El material que gana electrones queda negativo; el que los pierde,
    positivo, con carga de igual valor. Presenta la serie de la guía (vidrio, papel, lana,
    polietileno, PVC): al frotar dos materiales, el de más a la izquierda queda positivo. Clave:
    el signo depende del \emph{par} (la lana es $+$ frente al PVC y $-$ frente al papel), y
    las series de distintos laboratorios no coinciden del todo (humedad, superficies).\\
  Contacto & 35–55 & Péndulo eléctrico: una bolita de papel aluminio colgada de un hilo. El
    tubo de PVC frotado la atrae; al tocarla, la bolita sale repelida. Pregunta: ¿por qué ahora
    se repelen? (Quedó con carga del mismo signo que el tubo.) Por contacto los dos cuerpos
    quedan con el mismo signo.\\
  Conductores & 55–75 & Frota una varilla metálica sostenida con la mano: no atrae papelitos
    (la carga se va por el cuerpo); el PVC sí. Conductores (metales, agua con sales, el
    cuerpo humano), aislantes (vidrio, plástico, madera seca, caucho) y semiconductores
    (silicio). Lectura de la aplicación de la guía: el carrotanque conectado a tierra y la
    pulsera antiestática.\\
  Práctica & 75–92 & En parejas, primer punto de la Tarea 3 (la serie); se socializa el
    literal a).\\
  Cierre & 92–100 & Se asigna la Tarea 3 para el lunes 21. Pregunta para el lunes: si un
    cuerpo cargado se acerca a otro \emph{sin tocarlo}, ¿puede atraerlo?\\
  \bottomrule
\end{tabularx}}
\end{center}

\begin{nota}
Lleva el tubo de PVC y la lana, una bolsa de polietileno, un vaso de vidrio liso, hojas de
papel, papelitos, una bolita de papel aluminio colgada de un hilo de coser (con cinta a una
regla apoyada entre dos libros) y una varilla metálica o una cuchara. Con la humedad de Cali
conviene secar los materiales antes (secador de pelo o al sol).
\end{nota}

\end{document}
